\documentclass[12pt]{article}

\usepackage[margin=1in]{geometry}
\usepackage{amsmath,amssymb}
\usepackage{booktabs}
\usepackage{parskip}

\title{Wing Area to Wing Length Ratio}
\author{Semion D}
\date{\today}

\begin{document}

\maketitle

\section{Introduction} 

Let's say we have 2 paper airplanes, one of them is really long, but at the same time narrow, and on the other hand we have a short but wide airplane. Both of them have the same wing area. The question is: what will fly longer and over longer distances, and what will be more effective? This will be a problem to solve in this short essay. It will be divided into 3 main parts: Hypothesis, which will make a prediction of the outcome based on math, physics and logic; next is an experiment, which will check the hypothesis; and the last one is the conclusion, which will summarize everything into a simple logic-math statement.   

\[
***********************
\]

\newpage

\section{Hypothesis}

\subsection{Forces}

At first we need to define the forces that act upon a complex airplane-looking object. 

Here is the list:

\renewcommand{\arraystretch}{2}
\begin{center}
\begin{tabular}{c|l}
\toprule
\textbf{Force} & \textbf{Representation} \\
\midrule
\textbf{Gravity} & \textbf{Representation through constant $g = 9.81$ and mass of the object} \\
\midrule
\textbf{Lift} & \textbf{Representation in Newtons as a variable $L$ in formulas} \\
\midrule
\textbf{Drag} & \textbf{Representation in Newtons as a variable $D$ in formulas} \\
\midrule
\textbf{Thrust} & \textbf{Representation in Newtons as a variable $T$ in formulas} \\
\bottomrule
\end{tabular}
\end{center}

\vspace{5mm}

In this case we will not use thrust for several reasons: it does not affect the efficiency of the airplane itself, it just creates a continuous kinetic energy input; and it is harder to simulate on a small experimental scale. The only way where thrust will be included is the power input when we throw the airplane. So what's left is Gravity (or weight), Lift, and Drag.

Lift is a force that pushes the plane upwards. It happens because the area under the wing has higher air pressure than the top, and this pressure difference moves the airplane up. It is calculated through this formula:

\[
L = \frac{1}{2}\rho v^2 S C_L
\]

\renewcommand{\arraystretch}{2}
\begin{center}
\begin{tabular}{c|l|l}
\toprule
\textbf{Variable} & \textbf{Meaning} & \textbf{Unit} \\
\midrule
\textbf{$L$} & \textbf{Lift} & \textbf{Newtons ($N$)} \\
\midrule
\textbf{$\rho$} & \textbf{Air density} & $\mathrm{kg/m^3}$ \\
\midrule
\textbf{$v$} & \textbf{Air velocity} & $\mathrm{m/s}$ \\
\midrule
\textbf{$S$} & \textbf{Wing area} & $\mathrm{m^2}$ \\
\midrule
\textbf{$C_L$} & \textbf{Coefficient of lift} & \textbf{Scalar} \\
\bottomrule
\end{tabular}
\end{center}

\vspace{5mm}

$C_L$ is a scalar from 0 to 1 and combines in itself many wing parameters, one of the main of them is an angle of attack of the wing $\alfa$ and it represents the angle of the wing. Here is the formula for it:

\[
C_L = C_L0 + a \alfa
\]

$a$ is the lift-curve slope, and in general it tells us the speed at which $C_L$ grows when $\alfa$ grows. Basically it is the dirivitive of a linear function $C_L$. $C_L0$ is the lift coefficient at \(0^\circ\) angle of attack, which might be positive or negative, but in this case we will consider this equal to 0. \\

\subsection{Other Formulas}

Before actually calculating parameters of our 2 prototypes we need to establish the difference between them using math. Engineers did it already long time ago and came up with this fromula:

\[
AR = \frac{b^2}{S}
\]

$AR$ is the apect ratio, which represents the relationship between wing area $S$ and wing span $b$ which is a distance from one tip of the wing to another. This formula will help us to separate 2 ptototypes on a mathematical level. One more formula that will make our life easier is a formula for the lift-curve slope:

$$ a=\frac{a_0}{1+\frac{a_0}{\pi e AR}} $$

where:

\(a\) = lift-curve slope
\(a_0\) = 2D airfoil lift-curve slope
\(e\) = Oswald efficiency factor
\(AR\) = aspect ratio

$a_0$ tells you how quickly the lift coefficient \(C_L\) changes with angle of attack for an idealized wing cross-section extending infinitely in the spanwise direction (was not expecting that i will write smth so clever). But counting that wing is a finite region, air starts to behave differently, creating friction and drag, which is included and considered in $a$. For our project $a_0$ will be $ \boxed{a_0\approx2\pi\ \text{rad}^{-1}} $. 




\end{document}
